\section{The Field Nobody Approved}
\label{sec:ch17-the-field-nobody-approved}

\index{ownership!the word the toolchain does not have|(}%
Chapter~\ref{ch:breaking-changes} closed on a type three services declare and
none of them owns. I used the word \emph{owns} there the way everybody uses
it, as though it named something. It does not name anything the toolchain
holds. There is no owner in the composed schema, none in the router execution
config, none in the Apollo Federation directive set, and the word appears
nowhere in the Composite Schemas draft
chapter~\ref{ch:federation-model} decided against building on. What a team
means by it is real, and it lives entirely in the conversation they have with
each other.

\index{ownership!two edits to one document}%
So take two edits that a team could make on a Tuesday afternoon, both to the
Ratings document, both without opening either of the other three. The Ratings
service owns no session row; chapter~\ref{ch:second-third-subgraph} gave it
three fields on a \code{Session} it enters through a key. The four services
themselves are untouched in every run below, and only the document handed to
the composer differs.

\subsection{A Field on Somebody Else's Type}

\index{Session type@\code{Session} type!a fourth service field}%
The first edit adds a field. \code{flagged: Boolean}, on the Ratings copy of
\code{Session}, of the kind a moderation feature needs. Compose the four
documents and the composer answers as it answers everything in
chapter~\ref{ch:breaking-changes} that is not a routing problem:

\begin{minted}[fontsize=\footnotesize,breakanywhere]{text}
Router execution config successfully written to ".../router.json".
\end{minted}

\index{composition!accepts a field on a type the service owns no row of}%
Exit zero. The composed schema gains exactly two lines, the description and
the field, so \code{Session} in the graph every client reads now carries a
field the Sessions team never wrote:

\begin{minted}[fontsize=\footnotesize]{text}
206a207,208
>   """Whether this session was flagged by a moderator."""
>   flagged: Boolean
\end{minted}

\index{ownership!adding a field asks nobody}%
\code{diff} writes an addition as \code{a}, so that first line is saying two
lines arrived after line 206 and nothing else moved. The Sessions document was
read on the way, because composition reads all four, and it was not consulted.
Nothing in the merge asks whether the team that defines \code{Session} wanted
a field it never declared, because nothing in the merge has a way to ask. And
this half is not even a complaint: it is the property
chapter~\ref{ch:do-you-need-this} sold federation on. Your deployment, yours
alone, and a field you can add without a meeting.

\subsection{A Field That Was Already Somebody Else's}

\index{override@\code{"@override}!one line takes a field}%
The second edit is different in kind and identical in cost. \code{Session}
already has an \code{abstract}, declared in
chapter~\ref{ch:hot-chocolate-zero} and answered out of the Sessions service's
own table ever since. Declare it in
the Ratings document with one directive on it, \code{@override(from:
"sessions")}, which is Apollo's mechanism for moving a field from one subgraph
to another~\autocite{apollodirectives}. Compose the same four documents, the
Sessions one still byte for byte what the service exports:

\begin{minted}[fontsize=\footnotesize,breakanywhere]{text}
Router execution config successfully written to ".../router.json".
\end{minted}

\index{composition!a field changes hands at exit zero}%
Exit zero, and the same one line. What changed is inside the config, in the
list each subgraph carries of the fields it is recorded as able to answer for
\code{Session}:

\begin{center}
\footnotesize
\begin{tabular}{l l}
\toprule
Subgraph & Fields of \code{Session} it can answer \tabularnewline
\midrule
\code{sessions}, before & \code{id}, \code{speaker}, \code{title},
  \code{abstract}, \code{startsAt}, \code{durationMinutes} \tabularnewline
\code{sessions}, after & \code{id}, \code{speaker}, \code{title},
  \code{startsAt}, \code{durationMinutes} \tabularnewline
\addlinespace
\code{ratings}, before & \code{averageScore}, \code{ratingCount},
  \code{feedbackUrl}, \code{id} \tabularnewline
\code{ratings}, after & \code{abstract}, \code{averageScore},
  \code{ratingCount}, \code{feedbackUrl}, \code{id} \tabularnewline
\bottomrule
\end{tabular}
\end{center}

\index{supergraph!a change of owner does not reach it}%
\code{abstract} has left one service and arrived at another. And the schema a
client reads is byte for byte the one it was before the edit: 247 lines,
\code{diff} silent, every field in the same position with the same type and
the same description. A client cannot see this change. It cannot introspect it
either, which is what separates it from every change in
chapter~\ref{ch:breaking-changes}: there, the schema moved and no check said
so, and here nothing moves at all except a routing table nobody outside the
graph ever reads.

\subsection{What the Router Does With It}

\index{override@\code{"@override}!is a claim, not a transfer}%
\code{@override} is a claim. Composition takes the claim on trust: it never
asks the Ratings service whether it can produce an abstract. Nothing in
composition asks a service anything at all, which is what makes it runnable on
a laptop with the graph switched off and is not a defect. It does mean that
the first thing in the whole system to check the claim is a request.

\index{HTTP 200!a field with a new owner and no data behind it}%
Start a router on that config, with the same four unmodified services behind
it, and send the page query this book has printed since
chapter~\ref{ch:data-without-n-plus-1}, with one field added:

\begin{minted}{graphql}
{ sessions(first: 10) { nodes { title abstract } } }
\end{minted}

\index{DOWNSTREAM\_SERVICE\_ERROR@\code{DOWNSTREAM\_SERVICE\_ERROR}!a field the new owner never had}%
Here is the \code{errors} half of what comes back, with the \code{locations}
array and the \code{data} key left out to keep the two messages next to each
other. The outer one is the router's and the inner one is the Ratings service
answering it:

\begin{minted}[fontsize=\footnotesize,breakanywhere]{json}
{"errors":[{"message":"Failed to fetch from Subgraph 'ratings' at Path 'sessions.nodes'.","extensions":{"errors":[{"message":"The field `abstract` does not exist on the type `Session`.","extensions":{"code":"DOWNSTREAM_SERVICE_ERROR"}}],"serviceName":"ratings","statusCode":200}}]}
\end{minted}

\index{HTTP 200!the graph still answers the rest of the page}%
Two hundred, as everything in chapter~\ref{ch:breaking-changes} was. The
\code{data} key that block leaves out is intact: four sessions, every title
where it was, and \code{abstract} null on all four. One field of one type
stopped working, in a graph nobody deployed a line of code to, and the two
sentences that say so are one inside the other.

\index{ownership!nobody moved the data}%
Nobody moved the data. That is the whole of what happened: a directive said
the Ratings service resolves \code{abstract} now, four documents composed
around the claim, a router loaded the result, and the service behind it has
never heard of the field. In a real migration the resolver ships first and the
directive follows, and the reason that ordering is a convention rather than a
rule is that nothing between the edit and the request enforces it.

\index{composition!refused nothing in this section}%
Two edits, one document each, and the composer refused neither. It does refuse
something, and what it refuses is where this chapter turns.
\index{ownership!the word the toolchain does not have|)}%
